\SourceLeaf{349}%
die Natur der krummen Fläche durch eine Gleichung zwischen \(x,y,z\), bestimmt ist, und noch zusammengesetzter wird jener, wenn die Natur der krummen Fläche dadurch gegeben ist, dass \(x,y,z\) in der Gestalt von Functionen zweier neuen veränderlichen Grössen \(p,q\) dargestellt sind. Im letzten Fall enthält der Ausdruck funfzehn Elemente, nemlich die partiellen Differentialquotienten der ersten und zweiten Ordnung von \(x,y,z\) nach \(p\) und \(q\): allein er ist weniger wichtig an sich, als weil er den Übergang zu einem andern bahnt, der zu den merkwürdigsten Sätzen in dieser Lehre gerechnet werden muss. Bei jener Art, die Natur der krummen Fläche darzustellen, hat der allgemeine Ausdruck für irgend ein Linearelement auf derselben,
\[
\text{oder für}\quad\BareRoot(\mathrm{d}x^2+\mathrm{d}y^2+\mathrm{d}z^2),\quad\text{die Form}\quad\BareRoot(E\mathrm{d}x^2+2F\mathrm{d}x\mathord{.}\mathrm{d}y+G\mathrm{d}y^2)
\]
wo \(E,F.\ G\) wiederum Functionen von \(p\) und \(q\) werden, der erwähnte neue Ausdruck für das Krümmungsmaass enthält nun bloss diese Grössen, und ihre partiellen Differentialquotienten der ersten und zweiten Ordnung. Man sieht also, dass zur Bestimmung des Krümmungsmaasses bloss die Kenntniss des allgemeinen Ausdrucks eines Linearelements erforderlich ist, ohne dass es der Ausdrücke für die Coordinaten \(x,y,z\) selbst bedarf. Eine unmittelbare Folge davon ist der merkwürdige Lehrsatz: Wenn eine krumme Fläche, oder ein Stück derselben auf eine andere Fläche abgewickelt werden kann, so bleibt nach der Abwickelung das Krümmungsmaass in jedem Punkt ungeändert. Als specieller Fall folgt hieraus ferner: In einer krummen Fläche, die in eine Ebene abgewickelt werden kann, ist das Krümmungsmaass überall \(=0\). Man leitet daraus sofort die characteristische Gleichung der in eine Ebene abwickelungsfähigen Flächen ab, nemlich, in so fern \(z\) als Function von \(x\) und \(y\) betrachtet wird,
\[
\frac{\mathrm{dd}z}{\mathrm{d}x^2}\mathord{.}\frac{\mathrm{dd}z}{\mathrm{d}y^2}-(\frac{\mathrm{dd}z}{\mathrm{d}x\mathord{.}\mathrm{d}y})^2=0
\]
eine Gleichung, die zwar längst bekannt, aber nach des Verf. Urtheil bisher nicht mit der erforderlichen Strenge bewiesen war.

Diese Sätze führen dahin, die Theorie der krummen Flächen aus einem neuen Gesichtspunkte zu betrachten, wo sich der Untersuchung ein weites noch ganz unangebautes Feld öffnet. Wenn man die Flächen nicht als Grenzen von Körpern, sondern als Körper, deren eine Dimension verschwindet, und zugleich als biegsam, aber nicht als dehnbar betrachtet, so begreift man, dass zweierlei
